\section{Repositorio del analisis}

El anexo digital del presente trabajo reune el codigo, las tablas finales y las salidas principales empleadas en el analisis de desigualdades urbanas en la ciudad de Valencia.

La estructura del repositorio es la siguiente:

\begin{itemize}
    \item \texttt{01\_codigo/}: scripts finales utilizados para depurar, jerarquizar y analizar la informacion.
    \item \texttt{02\_datos\_finales/}: tablas finales empleadas en el analisis y diccionario de variables.
    \item \texttt{03\_resultados/}: informes de correlaciones, resumenes y salidas finales, incluyendo el mapa interactivo.
    \item \texttt{04\_material\_complementario/}: materiales adicionales no imprescindibles para la lectura del texto principal.
\end{itemize}

El repositorio asociado puede consultarse en:

\begin{center}
\texttt{https://github.com/USUARIO/REPOSITORIO}
\end{center}

Este repositorio se aporta con fines de transparencia, trazabilidad y reproducibilidad del proceso de analisis.
